\section{Deployment}\label{deployment}

The system can be brought up in two ways: \textbf{Docker Compose}, for local
development and for running the automated test suites of validation,
and \textbf{Kubernetes} for a production-like deployment.
Both are declarative: bringing the system up or down is a single command.

\subsection{Prerequisites}\label{deployment-prerequisites}

No Node.js version is declared in any \texttt{package.json} (no
\texttt{engines} field, no \texttt{.nvmrc}); the authoritative version is
fixed at the container level where every service and the
frontend build \texttt{FROM node:24-alpine}, so it runs, whether or not Node is installed on the host. 
Consequently, the
only host-level prerequisites are:

\begin{itemize}
  \item \textbf{Docker} and \textbf{Docker Compose} (v2 syntax), for the
  Compose path.
  \item \textbf{Minikube} and \textbf{kubectl}, for the Kubernetes path.
\end{itemize}

\subsection{Running with Docker Compose}\label{deployment-compose}

\begin{sloppypar}
\paragraph{Configuration.} The file groups its variables by concern:
shared credentials and infrastructure (\texttt{MONGO\_ROOT\_USER},
\texttt{MONGO\_ROOT\_PASSWORD}, \texttt{JWT\_SECRET}, \texttt{REDIS\_URL}),
one MongoDB connection string per service (\texttt{*\_MONGO\_URI}), and one
HTTP port per service. The chatbot's optional third-party integration
(Mistral AI, \cref{technological-details}) is configured separately via
\texttt{MISTRAL\_API\_KEY}/\texttt{MISTRAL\_MODEL}: if unset, the chatbot
falls back to template-based answers instead of failing.
\end{sloppypar}

\paragraph{Two Compose files, two purposes.} \texttt{docker-compose.dev.yml} is the
actively maintained development environment with the end-to-end
suite of validation targets: it adds the \texttt{order-service},
gives every MongoDB container a named volume and a health check, exposes
every database on a host port for inspection with a GUI client such as
MongoDB Compass, and additionally runs the observability stack.
Both manage the CAD context as a genuine three-member replica set initialised by a one-shot
\texttt{mongo-cad-init} job that waits for all three members to be healthy.

\begin{sloppypar}
\paragraph{Bringing it up.} From \texttt{kompozer/}:
\begin{quote}
\texttt{docker compose -f docker-compose.dev.yml up --build}
\end{quote}
Expected outcome: every container reports healthy, the CAD replica set
finishes electing a primary within a few seconds, and the API gateway
becomes reachable at \url{http://localhost:3000} (health check:
\texttt{GET /health} returns the aggregated status of
fault-tolerance). The frontend is served separately in this
environment (\texttt{npm run dev:frontend} from the workspace root, a Vite
dev server on \url{http://localhost:5173} proxying \texttt{/api} to the
gateway); the automated test suites can then be run as described in
\cref{validation}. Tearing down: \texttt{docker compose -f
docker-compose.dev.yml down -v} (the \texttt{-v} also removes the named
volumes for a full data reset).
\end{sloppypar}

\subsection{Deploying to Kubernetes}\label{deployment-kubernetes}

\paragraph{Layout.} All manifests live in \texttt{kompozer/k8s/}, numbered to
convey dependency order and aggregated by a single
\texttt{kustomization.yaml} (namespace \texttt{kompozer}).
The gateway and the frontend are the only two \texttt{NodePort} Services while
every other Service is a \texttt{ClusterIP}, unreachable from outside the
cluster, matching the single-entry-point design of \cref{architecture}.

\begin{sloppypar}
\paragraph{Bringing it up.} On a machine with Minikube:
\begin{quote}
\texttt{minikube start --driver=docker --memory=6144 --cpus=4}\\
\emph{(build every service image directly into Minikube's Docker store)}\\
\texttt{minikube image build ...} \emph{(once per service/frontend image)}\\
\texttt{kubectl apply -k kompozer/k8s}
\end{quote}
Images are built locally and loaded straight into Minikube rather than
pulled from a registry (every application \texttt{Deployment} sets
\texttt{imagePullPolicy: Never}); only the infrastructure images (MongoDB,
Redis, Loki, Promtail, Grafana) are pulled from public registries as usual.
Expected outcome, checked with \texttt{kubectl rollout status
\url{statefulset/mongo-cad}} and \texttt{kubectl wait
--for=condition=complete \url{job/mongo-cad-init}}: the CAD replica set is
initialised and every Deployment reaches its desired replica count; the
application is then reachable via \texttt{minikube service frontend -n
kompozer --url} (and, separately, \texttt{api-gateway}/\texttt{grafana} the
same way). Tearing down: \texttt{kubectl delete -k kompozer/k8s}.
\end{sloppypar}

\paragraph{Demonstrating the distributed-systems properties.} Because this
deployment is closer to production than Compose, it is also where the
fault-tolerance guarantees can be simulated against
a real node/process failure rather than a fake: killing the current CAD
primary (\texttt{kubectl delete pod mongo-cad-0}) shows automatic re-election
and writes resuming (NFR3); killing the CAD service pod (\texttt{kubectl
delete pod -l app=cad-service}) shows an open collaborative session recover
via checkpoint/replay after Kubernetes restarts it (NFR4).

\subsection{Observability}\label{deployment-observability}

Both environments can additionally run a log-aggregation stack (
\textbf{Grafana}, \textbf{Loki} and \textbf{Promtail}) scoped
to development/troubleshooting rather than production use (in-memory index,
filesystem storage, anonymous admin access). Promtail runs once per Docker
host in Compose, and as a Kubernetes
\texttt{DaemonSet} in the cluster deployment, shipping every
container's logs to Loki; Grafana is pre-provisioned with Loki as a
datasource, so logs from any of the nine backend components can be filtered
and correlated from a single dashboard during manual testing.

